\setcounterpageref{ex}{0}
\TaoPhieuKiemTra{17}

\setcounter{ex}{0}
\setcounter{bt}{0}

\textbf{\large I. PHẦN TRẮC NGHIỆM (3,0 điểm)}

{\itshape Khoanh tròn vào chữ cái đặt trước câu trả lời đúng. Mỗi câu đúng được 0,5 điểm.}

\begin{ex} % Câu 1
Từ hai chữ số $4$ và $7$ lập được tất cả bao nhiêu số có hai chữ số khác nhau?
\choice
{$1$ số}
{\True $2$ số ($47$ và $74$)}
{$3$ số}
{$4$ số}
\loigiai{
Các số lập được là $47$ và $74$ (có $2$ số).\\
\textbf{Đáp án đúng: B}
}
\end{ex}

\begin{ex} % Câu 2
Tổng của số lớn nhất có một chữ số ($9$) và số bé nhất có hai chữ số ($10$) là
\choice
{$18$}
{$20$}
{$99$}
{\True $19$}
\loigiai{
Tổng: $9 + 10 = 19$.\\
\textbf{Đáp án đúng: D}
}
\end{ex}

\begin{ex} % Câu 3
Tính: $17\text{ }l - 9\text{ }l = \dots$
\choice
{\True $8\text{ }l$}
{$7\text{ }l$}
{$9\text{ }l$}
{$6\text{ }l$}
\loigiai{
$17 - 9 = 8\text{ }l$.\\
\textbf{Đáp án đúng: A}
}
\end{ex}

\begin{ex} % Câu 4
Hình tứ giác có tất cả bao nhiêu đỉnh?
\choice
{$3$ đỉnh}
{$5$ đỉnh}
{\True $4$ đỉnh}
{$6$ đỉnh}
\loigiai{
Hình tứ giác có $4$ đỉnh.\\
\textbf{Đáp án đúng: C}
}
\end{ex}

\begin{ex} % Câu 5
So sánh: $15 - 8 \dots 16 - 9$.
\choice
{$>$}
{\True $=$}
{$<$}
{Không so sánh được}
\loigiai{
$15 - 8 = 7$; $16 - 9 = 7$. Vì $7 = 7$ nên điền dấu $=$.\\
\textbf{Đáp án đúng: B}
}
\end{ex}

\begin{ex} % Câu 6
Tính: $8 + 7 - 6 = \dots$
\choice
{\True $9$}
{$8$}
{$10$}
{$7$}
\loigiai{
$8 + 7 = 15$; $15 - 6 = 9$.\\
\textbf{Đáp án đúng: A}
}
\end{ex}

\vspace{4mm}
\textbf{\large II. PHẦN TỰ LUẬN (7,0 điểm)}

{\itshape Thí sinh trình bày chi tiết bài làm vào giấy kiểm tra.}

\begin{bt}[1,0 điểm]
Tính nhẩm:
\begin{multicols}{2}
\begin{enumerate}[a)]
  \item $9 + 6$
  \item $7 + 7$
  \item $14 - 6$
  \item $16 - 7$
\end{enumerate}
\end{multicols}
\loigiai{
\begin{itemize}
  \item $9 + 6 = 15$
  \item $7 + 7 = 14$
  \item $14 - 6 = 8$
  \item $16 - 7 = 9$
\end{itemize}
}
\end{bt}

\begin{bt}[1,5 điểm]
Đặt tính rồi tính:
\begin{multicols}{3}
\begin{enumerate}[a)]
  \item $27 + 33$
  \item $89 - 32$
  \item $50 + 42$
\end{enumerate}
\end{multicols}
\loigiai{
\begin{itemize}
  \item $27 + 33 = 60$
  \item $89 - 32 = 57$
  \item $50 + 42 = 92$
\end{itemize}
}
\end{bt}

\begin{bt}[1,0 điểm]
Tính giá trị biểu thức kèm đơn vị đo:
\begin{enumerate}[a)]
  \item $33\text{ kg} + 8\text{ kg} - 11\text{ kg}$
  \item $42\text{ }l - 13\text{ }l + 7\text{ }l$
\end{enumerate}
\loigiai{
\begin{enumerate}[a)]
  \item $33\text{ kg} + 8\text{ kg} - 11\text{ kg} = 41\text{ kg} - 11\text{ kg} = 30\text{ kg}$.
  \item $42\text{ }l - 13\text{ }l + 7\text{ }l = 29\text{ }l + 7\text{ }l = 36\text{ }l$.
\end{enumerate}
}
\end{bt}

\begin{bt}[1,5 điểm]
Bạn Minh có $32$ bức tranh. Bạn Phong có nhiều hơn bạn Minh $9$ bức tranh. Hỏi bạn Phong có bao nhiêu bức tranh?
\loigiai{
Số bức tranh bạn Phong có là:
\[ 32 + 9 = 41 \text{ (bức tranh)} \]
\textbf{Đáp số:} $41$ bức tranh.
}
\end{bt}

\begin{bt}[1,0 điểm]
Một thùng chứa đầy $52\text{ }l$ nước mắm. Buổi sáng cửa hàng đã bán ra $13\text{ }l$. Hỏi trong thùng còn lại bao nhiêu lít nước mắm?
\loigiai{
Trong thùng còn lại số lít nước mắm là:
\[ 52 - 13 = 39\text{ }l \]
\textbf{Đáp số:} $39\text{ }l$ nước mắm.
}
\end{bt}

\begin{bt}[1,0 điểm]
Tìm số có hai chữ số, biết chữ số hàng chục là số lớn nhất có một chữ số, chữ số hàng đơn vị kém chữ số hàng chục $4$ đơn vị.
\loigiai{
Chữ số hàng chục là $9$.\\
Chữ số hàng đơn vị là: $9 - 4 = 5$.\\
Vậy số cần tìm là $95$.\\
\textbf{Đáp số:} $95$.
}
\end{bt}
